\def\velkostpisma{9}
\def\velkostriadku{12}
\input makra.tex % nacitanie Ivanom pripravenych nastaveni a prikazov
\hyphenation{star-šov-stvo} % rozdelenie slov na konci riadku, treba tu uviest slova, ktore sam nepozna

\spravodaj{10}{2026}


\clanok{{\it ČO SA DEJE V~OKTÓBRI}}


\clanok{Verš na mesiac}
V~októbri sa učíme nový verš, ktorý dostali pre rok~2026 naši diakoni:

„Hospodinovo oko spočíva na tých, čo sa Ho boja, na tých, čo očakávajú Jeho milosť.“ (Ž~33,18)


\clanok{Uvítacia služba}
Od septembra sme obnovili uvítaciu službu pred bohoslužbami, ktorá vytvára priestor pre rozhovory a lepšie vzájomné spoznávanie sa.
Chceme tak prispieť k~tomu, aby sa u~nás každý cítil prijatý a ako doma. Veríme, že takáto služba pomôže nielen novým tváram medzi nami, ale posilní aj vzťahy v~celom našom spoločenstve.

Vytvorili sme možnosť prihlásiť sa na službu v~online tabuľke.
Ospravedlňujeme sa za to, že prvé týždne nefungovala ako by mala. Chybu sme odstránili, táto možnosť je k~dispozícii. Akokoľvek, ak by ste mali záujem zapojiť sa, ozvite sa Filipovi Barkóczimu.


\clanok{Biblické štúdium v~anglickom jazyku}
Plettovci pozývajú na biblické štúdium v~angličtine.
Začiatok je v~stredu 7.~10.~2026.
Bližšie informácie Vám poskytne Daniel.


\clanok{Rozhovory na Palisádach -- sobota 10.~10.~2026}
Keď sa povie paliatívna starostlivosť, veľa ľudí si predstaví koniec.
Táto starostlivosť však nie je o~pridávaní dní do života, ale pridávaní života do dní, ktoré zostávajú. Príďte si vypočuť prednášku na túto tému a položiť otázky našej hosťke MUDr. Andrei Škripekovej, PhD. v~sobotu 10.~10.~2026
o~9.30~hod. v~našej modlitebni na Palisádach 27/A~v~Bratislave.
Všetci ste srdečne vítaní aj s~priateľmi, rodinou a všetkými, ktorým táto téma prinesie úžitok.


\clanok{Oblastné vďakyvzdanie -- nedeľa 11.~10.}
Aj tento rok organizuje náš zbor oblastné vďakyvzdanie. Uskutoční sa v~SÚZA (Správa účelových zariadení) na Drotárskej ceste č.~4748.
Program začína o~9.30~hod., v~tomto čase sa bohoslužby v~našej modlitebni na Palisádach neuskutočnia.
Celý program  je zverejnený na plagáte v~modlitebni a rovnako aj v~e-mailových oznamoch aj s~organizačnými pokynmi.
Ak viete o~niektorých, ktorí by sa nevedeli dostať k~týmto informáciam, prosíme, pomôžte im v~tom.

Tešíme sa na spoločný čas s~bratmi, sestrami, priateľmi a deťmi zo zborov našej oblasti.
Budeme spoločne vzdávať vďaku dobrému Bohu za dobré dary podľa toho, ako Pavol píše Kolosenským: „Pevne sa držte viery, ako vám ju vštepili, a denne ďakujte Bohu za všetko, čím Vás obdaril.“ (Kol~2,7)


\clanok{Deň diakonie -- nedeľa 18.~10.}
V~nedeľu 18.~10.~2026 si pripomenieme Deň diakonie. Pre tento účel zverejňujeme
(čítajte dolu nižšie) list od koordinátorov Dobrodiny, v~ktorom nás prosia
o~zorganizovanie finančnej zbierky pre tento deň.
Zbierka sa uskutoční v~našom zbore 18.~10.~2026.


\clanok{Stretnutie sestier -- streda 21.~10.}
Sestry sa stretnú v~stredu 21.~10.~2026 o~17.30~hod. na Zrínskeho~2.
Témou stretnutia bude „Neotrasiteľné kráľovstvo“.
Je to téma z~konferencie EBWU (Únia baptistických žien Európy), ktorú pripraví sestra Jarka Cihová.

„A~tak keďže prijímame neotrasiteľné kráľovstvo, zachovávajme si milosť a pomocou nej slúžme Bohu s~úctou a bázňou, aby sme sa mu páčili.“ (Žid~12,28)
Všetky sestry sú srdečne vítané.


\clanok{Zborové členské zhromaždenie -- nedeľa 25.~10.}
V~nedeľu 25.~10.~2026 sa uskutoční zborové členské zhromaždenie. Rezervujte si, prosím, tento termín. Čas oznámime neskôr a program sa predbežne bude týkať DKDZ a smerovania F.~Barkócziho v~našom zbore. Účať členov je veľmi vítaná.


\clanok{Obnova fasády na Palisádach}
V~rámci programu Ministerstva kultúry SR „Obnovme si svoj dom“ sme žiadali
o~finančný príspevok na obnovu fasády nášho kostola. Žiaľ, naša žiadosť nebola podporená. Grant sme však považovali len za pomoc nášmu rozpočtu, a preto práce pokračujú podľa plánu. Ďakujeme vám za podporu a modlitby.


\clanok{{\it MIMOZBOROVÉ POZVANIA}}


\clanok{Storočnica v~Miloslavovskom zbore}
V~tomto roku prežívame v~Miloslavove veľkú vďačnosť za Božiu vernosť -- pripomíname si vzácne jubileum 100.~výročia otvorenia našej modlitebne a osamostatnenia cirkevného zboru.

Za týmto storočím vidíme nespočetné množstvo príbehov Božej milosti, služby a spoločenstva, ktoré formovali nielen náš zbor, ale aj celú obec. Veľmi si vážime naše bratské vzťahy v~rámci BJB a bola by pre nás veľká radosť a povzbudenie, keby ste toto výročie oslávili spolu s~nami.

Srdečne Vás pozývame na slávnostné bohoslužby, ktoré sa uskutočnia v~nedeľu 18.~10.~2026 o~9.45~hod. v~modlitebni BJB na Centrálnej ulici~140 v~Miloslavove. Bude pre nás poctou, ak sa zástupcovia Vášho zboru alebo kazatelia pri tejto príležitosti prihovoria zhromaždeniu s~krátkym bratským pozdravom.

Okrem hlavných bohoslužieb Vás i členov Vašich zborov pozývame aj na sprievodné víkendové podujatia:
\begitems
* sobota 17.~10. o~17.00~hod. -- Koncert spevokolu JAS
* nedeľa 18.~10. o~15.00~hod. -- Koncert modlitieb a chvál MS Praise Band
\enditems

Z~dôvodu organizačného zabezpečenia Vás prosíme o~krátke potvrdenie účasti
(do 10.~10.~2026 e-mailom: \email{miroslavt@gmail.com} alebo telefonicky 0902~677~845)

\autor{S~láskou a bratským pozdravom v~Kristovi,}
\autor{Mgr. Miroslav Tóth, kazateľ CZ BJB Miloslavov}


\clanok{Konferencia Aglow}
Sestra Rút Krajčíová pozýva sestry na konferenciu Aglow Slovakia.
Uskutoční sa od 23.~10. do 25.~10.~2026 v~Poprade. Viac informácií je v~pozvánke v~e-maily alebo priamo u~sestry Rút.


\clanok{{\it SPRÁVY ZO SEPTEMBRA}}


\clanok{Správy zo staršovstva}
Od začiatku prázdnin sme mali štyri stretnutia.
Na augustovom stretnutí sme hovorili o~nominácii brata kazateľa Janka Szőllősa na predsedu Rady BJB na Slovensku. Vidíme potrebu podporiť túto službu pre celú jednotu a predo sme odporučili brata Janka Szőllősa za predsedu, prípadne podpredsedu Rady BJB na Slovensku.
Ďalej sme hovorili o~správe budovy na Súľovskej.

Br.~kazateľ Dávid M.~Chuchút nám prezentoval svoju predstavu tematického zamerania nedeľných bohoslužieb, program kázaní aj zmenu organizácie spoločných bohoslužieb s~nadväznosťou na biblické hodiny. Podľa jeho návrhu preberáme od septembra 1.~list ap.~Petra. Organizáciu bohoslužieb sme upravili tak, aby hlavná časť nedeľných bohoslužieb bola zameraná na Božiu oslavu a vstupy zamerané na zborový život sme sústredili ako záverečnú časť našich stretnutí.
Pripravovali sme inštaláciu br.~D.~M.~Chuchúta za kazateľa zboru, ktorú sme mali 6.~9. za účasti predsedu Rady BJB br.~Benjamína Uhrina, Dávidovej manželky a rodiny.

Na ďalších stretnutiach v~septembri sme hovorili o~príprave vďakyvzdania, ktoré chceme mať 11.~10. s~bratmi a sestrami zo zborov z~našej oblasti.
Tiež sme sa pripravovali na stretnutia KDZ, ktoré budú v~októbri a novembri a na ktorých prebehnú aj spomínané voľby do Rady BJB.

S~KDZ súvisí aj potreba mať ZČZ, na ktorom budeme informovať o~témach z~KDZ a prípadne hlasovať k~prerokovávaným témam. ZČZ plánujeme na nedeľu 25.~10.

Stretnutia staršovstva sa snažíme zmeniť zo zamerania na riešenie technických otázok zboru a viac sa zameriavať na duchovné potreby členov a priateľov zboru.
Preto prosím o~vaše modlitby za múdrosť a Božie vedenie vo všetkom, čo dáva do našich osobných životov aj do života nášho spoločenstva.
S~láskou a modlitbou
\autor{Peter Pribula}


\clanok{Dobrodina}
Milé sestry, milí bratia, v~nedeľu 18.~10. si spoločne pripomenieme Deň diakonie. Je to príležitosť poďakovať za službu, ktorou si navzájom pomáhame, a pamätať na ľudí, ktorí potrebujú našu blízkosť a starostlivosť. Pri tejto príležitosti sa s~vami chceme podeliť o~radostné novinky z~Dobrodiny a pozvať vás, aby ste boli aj naďalej jej súčasťou.

Dobrodina je projektom Bratskej jednoty baptistov na Slovensku. V~Bernolákove postupne pripravujeme miesto, kde sa prepoja služby pre seniorov s~materskou školou, komunitným bývaním a zdravotným strediskom. Túžime, aby tu ľudia rôznych generácií nachádzali pomoc, prijatie a blízkosť a aby nikto nemusel zostať so svojou krehkosťou, zraniteľnosťou či ťažkosťami sám.

Nadväzujeme tak na službu našich bratov a sestier, ktorí už pred viac ako sto rokmi premieňali svoju vieru na konkrétnu pomoc a skutky lásky. S~vďačnosťou na ich službu spomíname a chceme v~nej pokračovať spôsobom, ktorý odpovedá na dnešné potreby.

Od minulého Dňa diakonie sme sa posunuli o~veľký krok bližšie k~uskutočneniu tejto vízie. Po niekoľkých rokoch príprav sme na konci septembra podpísali
s~Ministerstvom práce, sociálnych vecí a rodiny SR zmluvu o~nenávratnom finančnom príspevku vo výške viac ako 3,3~milióna~€ na vybudovanie špecializovaného zariadenia.

Máme z~toho veľkú radosť. Vďaka tejto podpore plánujeme už v~roku~2028 otvoriť prvú časť Dobrodiny -- domov so špecializovanou starostlivosťou pre 30~ľudí, ktorí sú v~dôsledku Alzheimerovej či Parkinsonovej choroby alebo iných neurodegeneratívnych ochorení odkázaní na pomoc druhých. Chceme im ponúknuť dôstojnú a láskavú starostlivosť a ich rodinám istotu, že ich blízki sú
v~dobrých rukách.

Za získaný príspevok z~Ministerstva sme veľmi vďační. Nepokrýva však všetky náklady na vybudovanie zariadenia a podmienkou jeho poskytnutia je aj naša finančná spoluúčasť -- spolufinancovanie projektu z~vlastných zdrojov. Aj preto potrebujeme pomoc cirkevných zborov, jednotlivcov a ďalších partnerov. Vaše dary nám pomôžu zabezpečiť tento vlastný podiel a doplniť získanú podporu, aby sme mohli Dobrodinu spoločne priviesť k~životu a otvoriť jej dvere tým, ktorí potrebujú starostlivosť.

Chceli by sme vás preto poprosiť, aby ste pri príležitosti Dňa diakonie podľa svojich možností usporiadali vo vašom zbore finančnú zbierku na podporu Dobrodiny. Budeme vďační za každý príspevok, ktorým sa stanete súčasťou tejto spoločnej služby. Výnos zbierky môžete poslať na účet Účelového zariadenia RODINA: SK62 0900 0000 0051 7799 3135.

Prosíme vás aj o~modlitby za ďalšie kroky, za ľudí, ktorí sa na príprave Dobrodiny podieľajú, za budúcich zamestnancov i za tých, ktorým bude raz slúžiť. Pri pohľade na doterajšiu cestu Dobrodiny si s~vďačnosťou uvedomujeme, koľko lásky, ochoty a vytrvalosti do nej už mnohí vložili. Povzbudením sú nám slová apoštola Pavla: „Vždy, keď sa modlíme, ďakujeme za vás Bohu, Otcovi nášho Pána Ježiša Krista. Veď sme počuli o~vašej viere v~Ježiša Krista a o~láske, ktorú máte k~všetkým svätým pre nádej, ktorá je vám pripravená v~nebesiach a
o~ktorej ste už skôr počuli v~pravdivom slove evanjelia, ktoré k~vám prišlo. Rovnako ako na celom svete aj u~vás prináša ovocie a rastie od toho dňa, keď ste počuli a poznali Božiu milosť v~pravde.“ (Kol~1,3-6)

Ďakujeme vám za vaše modlitby, štedrosť a mnohé povzbudzujúce slová. Uvedomujeme si, že aj vo svojich zboroch nesiete mnohé potreby a verne slúžite ľuďom okolo vás. O~to viac si vážime, že pamätáte aj na Dobrodinu. Prajeme vám, aby ste vo svojej službe zažívali Božiu blízkosť, silu a radosť.

\autor{S~vďačnosťou a prianím Božieho pokoja,}
\autor{Lea Miklošová, riaditeľka Účelového zariadenia RODINA a}
\autor{Janine Lacho, projektová koordinátorka rozvoja Dobrodiny.}
\autor{Web: \ulink[https://www.dobrodina.sk]{www.dobrodina.sk}.}


\clanok{Zbierky v~septembri}
\table{lr}{
Na misiu           & 389~€\cr
Na investičný fond &  95~€\cr
}
\vskip1em

Pripomíname, že zbierky sa u~nás konajú pravidelne takto:
\vskip-1ex\begitems
* každú 2.~nedeľu v~mesiaci je zbierka venovaná misii a
* každú 4.~nedeľu je zbierka na investície.
\enditems

Prispieť je možné formou zbierky počas 4.~nedieľ v~mesiaci alebo bankovým prevodom na zborový účet, prípadne aj osobne s.~Ľubke Kohútovej.

Bankové spojenie: SK36 0900 0000 0000 1147 1836, SWIFT: GIBASKBX



\n 2.	10.	Peter	ANTALÍK;
\n 6.	10.	Daniel	BALÁŽ;
\n 12.	10.	Barbora	PRIBULOVÁ;
\n 14.	10.	Martin	SIMON;
\n 20.	10.	Ida	PUČEKOVÁ;
\n 22.	10.	Hana	HALAMIČKOVÁ;
\n 25.	10.	Vladimír	IRA;
\n 26.	10.	Martin	HOVORKA;
\n 27.	10.	Miriam	KRÁĽOVÁ;
\n 28.	10.	František	VRABČEK;
\n 28.	10.	Ľubomír	SYČ;
\n 29.	10.	Zuzana PAŘÍZKOVÁ;
\narodeniny


\program{
\p  1 ; št ; 18.00 ; Biblická hodina (D.~M.~Chuchút) ;.;;
\p  2 ; pi ; 17.30 ; Dorast ;.;;
\p  3 ; so ; 18.00 ; Mládež ;.;;
\p  4 ; ne ;  9.30 ; Bohoslužby (D.~M.~Chuchút + VP);.;;
\p  5 ; po ; 18.00 ; Modlitebná skupinka ;.;;
\p  6 ; ut ; 15.00 ; Biblická hodina pre seniorov (P.~Pivka) ;.;;
\p  7 ; st ;.;;.;;
\p  8 ; št ; 18.00 ; Biblická hodina (D.~M.~Chuchút) ;.;;
\p  9 ; pi ; 17.30 ; Dorast ;.;;
\p 10 ; so ; 18.00 ; Mládež ;.;;
\p 11 ; ne ;  9.30 ; Oblastné vďakyvzdanie v~SÚZA -- Bohoslužby na Palisádach NEBUDÚ  ;.;;
\p 12 ; po ; 18.00 ; Modlitebná skupinka / skupinka „Rastieme v~Kristovi“ ;.;;
\p 13 ; ut ; 15.00 ; Biblická hodina pre seniorov (P.~Pivka) ;.;;
\p 14 ; st ;.;;.;;
\p 15 ; št ; 18.00 ; Biblická hodina (D.~M.~Chuchút) ;.;;
\p 16 ; pi ; 17.30 ; Dorast ;.;;
\p 17 ; so ; 18.00 ; Mládež ;.;;
\p 18 ; ne ;  9.30 ; Bohoslužby (R.~Nemec) ;.;;
\p 19 ; po ; 18.00 ; Modlitebná skupinka ;.;;
\p 20 ; ut ; 15.00 ; Biblická hodina pre seniorov (P.~Pivka) ;.;;
\p 21 ; st ; 17.30 ; Stretnutie sestier ;.;;
\p 22 ; št ; 18.00 ; Biblická hodina (D.~M.~Chuchút) ;.;;
\p 23 ; pi ; 17.30 ; Dorast ;.;;
\p 24 ; so ; 18.00 ; Mládež ;.;;
\p 25 ; ne ;  9.30 ; Bohoslužby (D.~M.~Chuchút);.;;
\p 26 ; po ; 18.00 ; Modlitebná skupinka / skupinka „Rastieme v~Kristovi“ ;.;;
\p 27 ; ut ; 15.00 ; Biblická hodina pre seniorov (P.~Pivka) ;.;;
\p 28 ; st ;.;;.;;
\p 29 ; št ; 18.00 ; Biblická hodina (D.~M.~Chuchút) ;.;;
\p 30 ; pi ; 17.30 ; Dorast ;.;;
\p 31 ; so ; 18.00 ; Mládež ;.;;
}

\tiraz
\bye
